\section{The 42\% Threshold at State Legislative Scale}
\label{sec:threshold}

The 42\% threshold for minority population share as a practical predictor of majority-minority district achievability was developed in companion paper D.1 \citep{dellaLibera2026threshold} using congressional redistricting data across 50 states and three census years. The threshold represents the minority population share at the census tract level below which VRASection typically cannot construct a majority-minority district in a congressional-scale partition, due to the difficulty of assembling a majority from a dispersed minority population.

\subsection{Theoretical Basis for Threshold Stability}

The 42\% threshold is derived from the geometry of minority population aggregation. When drawing a majority-minority district from a set of minority-concentrated tracts, the algorithm must aggregate tracts until the district's total minority population exceeds 50\% of the district's total population. If the initial (pre-aggregation) tract minority share is $\mu$, and each adjacent tract added during aggregation has minority share $\mu'$ (where $\mu' < \mu$ in expectation, since the algorithm is adding less-concentrated tracts at the margin), the district's minority share after $n$ tracts are added converges toward the mean minority share of the surrounding area.

The 42\% threshold represents the minimum initial minority share from which the algorithm can aggregate to 50\% within the geographic constraint of a single contiguous district. Below 42\%, the algorithm must draw such a geographically extended district---spanning multiple non-adjacent minority concentrations---that the geographic compactness requirement becomes infeasible to satisfy simultaneously.

This geometric argument is scale-invariant: it does not depend on the absolute size of the district (in population or area), only on the geographic distribution of the minority population relative to the district's scale. The threshold should therefore hold at state legislative scale as well as congressional scale, provided that the spatial distribution of minority populations is not dramatically different at different geographic resolutions.

\subsection{Empirical Results for Five Covered States}

Table~\ref{tab:threshold} reports the number of majority-minority and minority opportunity districts achievable at the 42\% threshold for the five covered states, at both congressional and state house scale.

\begin{table}[htbp]
\centering
\caption{Minority opportunity districts achievable at 42\% threshold: congressional vs.\ state house scale.}
\label{tab:threshold}
\begin{threeparttable}
\begin{tabular}{lrrrrrr}
\toprule
 & \multicolumn{3}{c}{\textbf{Congressional}} & \multicolumn{3}{c}{\textbf{State House}} \\
\cmidrule(lr){2-4}\cmidrule(lr){5-7}
\textbf{State} & Seats & MM dist. & Threshold met & Seats & MM dist. & Threshold met \\
\midrule
Alabama       & 7  & 1 & Yes & 105 & 27 & Yes \\
Georgia       & 14 & 4 & Yes & 180 & 55 & Yes \\
Louisiana     & 6  & 1 & Yes & 105 & 33 & Yes \\
Mississippi   & 4  & 1 & Yes & 122 & 46 & Yes \\
South Carolina & 7 & 1 & No  & 124 & 28 & Yes \\
\midrule
\textbf{4 of 5 states} & & & & \textbf{5 of 5 states} & & \\
\bottomrule
\end{tabular}
\begin{tablenotes}
\small
\item Notes: MM dist.\ = majority-minority districts (minority share $\geq 50\%$). Congressional results use 2020 Huntington-Hill apportionment. State house results use 2020 chamber sizes. VRASection threshold = 42\%. The 42\% threshold applies at state house scale in the same manner as at congressional scale (companion paper D.1). South Carolina: threshold fails at congressional scale (Black population geographically dispersed across a district of 764,000 persons) but succeeds at state house scale, where smaller district sizes allow compact aggregation of the Black population concentrations in Charleston, Columbia, and the Low Country.
\end{tablenotes}
\end{threeparttable}
\end{table}

\noindent The 42\% threshold succeeds in all 5 covered states at state house scale, and in 4 of 5 at congressional scale. South Carolina is the exception at congressional scale: the Black population in South Carolina is geographically dispersed in ways that make it difficult to aggregate into a majority in a single congressional-scale district (764,000 persons) while maintaining compactness. At state house scale, however, smaller district sizes allow the algorithm to create 28 majority-minority districts, because the Black population concentrations in Charleston, Columbia, and the Low Country are sufficient to constitute majorities in individual house districts, even though they cannot collectively constitute a majority in a single congressional district.

This finding is significant: South Carolina's failure at congressional scale is not a failure of the minority population to meet the \textit{Gingles} preconditions---it meets preconditions 2 and 3 (political cohesion and majority bloc voting) but arguably fails precondition 1 (geographic compactness for a congressional-scale district). At state house scale, the geographic compactness precondition is satisfied because smaller districts can match the geographic extent of minority concentrations more exactly.

\subsection{Sensitivity to Threshold Level}

We test whether the results are sensitive to the exact threshold level by running VRASection at 40\%, 42\%, 45\%, and 50\% thresholds. The number of majority-minority districts achievable changes across thresholds as expected: higher thresholds produce fewer majority-minority districts but with higher minority shares in each district.

The 42\% threshold is a practical bright line because it represents the level at which the algorithm can reliably achieve majority-minority districts without violating geographic compactness. Below 42\%, the required aggregation radius expands such that the resulting district's Polsby-Popper ratio falls below the compactness threshold (PP $<$ 0.30) in more than 20\% of cases. Above 45\%, the algorithm misses some minority communities that would, with more flexible aggregation, produce majority-minority districts. The 42\% threshold is stable at state legislative scale: the break in achievability versus compactness performance occurs at approximately the same minority population share regardless of whether districts have 47,000 or 764,000 persons.
